\documentclass[10pt,a4paper]{amsart}
\usepackage[margin=19mm]{geometry}
\usepackage{amsmath,amssymb,mathtools}
\usepackage[scr=rsfs,cal=euler]{mathalfa}
\usepackage{xfrac}
\usepackage[unicode,hidelinks,bookmarks=false]{hyperref}
\newcommand{\ie}{i.e.\ }
\newcommand{\cf}{cf.\ }
\newcommand{\resp}{respectively\ }
\newcommand{\Subsection}{\subsection}
\providecommand{\nobreakdash}{\nobreak\hbox{-}}
\setlength{\emergencystretch}{2em}
\multlinegap0pt
\usepackage[shortlabels]{enumitem}
\usepackage{mathtools}
\usepackage{amssymb}
\usepackage{xcolor}
\usepackage{bm}
\usepackage{float}
\newtheorem{lemma}{Lemma}
\def\C{\mathbb{C}}
\def\P{\mathcal{P}}
\def\R{\mathbb{R}}
\renewcommand{\d}{\text{\rm d}}
\newcommand{\dD}{\partial\mathbb{D}}
\title{Sharp sign uncertainty for trigonometric polynomials}
\author{Tolibjon Ismoilov}
\date{Source: arXiv:2606.02299v1 (CC BY 4.0)}
\begin{document}
\maketitle

\noindent\emph{Source and licence.} Sharp sign uncertainty for trigonometric polynomials. Version arXiv:2606.02299v1 (CC BY 4.0), distributed by arXiv under the Creative Commons Attribution 4.0 International licence, http://creativecommons.org/licenses/by/4.0/, as recorded in the arXiv licence metadata for this exact version and shown on the abstract page https://arxiv.org/abs/2606.02299v1. Full licence notice: \texttt{licenses/arxiv-2606.02299-CC-BY-4.0.txt}.\par\smallskip

\noindent\textit{Editorial scope.} Counted unit: the article's lemma lem:quadrature-formula-for-laurent-polynomials (source0.tex lines 911--917). The article's complete original proof of it is delivered with it (source0.tex lines 919--955, one \texttt{proof} environment of 37 lines). No mathematics is written, repaired or re-derived: the statement and the proof are the article's own lines, in the article's own notation, and no retained line is substituted or re-flowed.  Retained uncounted as the source-local context the proof needs: lines 716--718; lines 834--848; lines 868--870; lines 873--875; lines 906--908.  Nothing was omitted from any retained span, so the package carries no omission record.  No retained line contains a citation, so the wrapper carries no bibliography and no bibliography entry.  No retained line contains a figure environment, an image-inclusion command or drawing code, so no image is delivered and the package declares no asset dependency.  Editorial formatting only, in addition: an AMS \texttt{amsart} preamble that restates the article's own declarations and macros used by the retained lines, namely the environments \texttt{lemma} on separate counters, together with the standard packages those lines need.  The retained environments print in this document as Lemma 1 in order of appearance, whereas the article prints the same items with the numbering of its own full document.  The complete native source of the article as distributed in the arXiv e-print for this exact version is bundled unchanged as \texttt{sources/201893/source0.tex}.
\begin{equation}\label{eq:reproducing-kernel-with-A-and-B}
    K(w, z)=\frac{2}{i}\frac{B(z)\overline{A(w)}-A(z)\overline{B(w)}}{1-\overline{w}z}, \quad \overline{w}z\neq 1.
\end{equation}

\begin{lemma}\label{lem:properties of orthogonal polynomials}
    Let $\mu$ be a probability measure on $\dD$.
    \begin{enumerate}[label=\textnormal{(\roman*)}]
        \item If $\d\mu(z)=\d\mu(\overline{z})$ holds,  then $\Phi_n(z)$ has real coefficients.
        \item $\varphi_n(z)$ has all its zeros in $\mathbb{D}$ and $\varphi^*_n(z)$ has all its zeros in $\C\setminus\overline{\mathbb{D}}$ for $n \geq 1$. Therefore, 
        \begin{equation*}
            \lvert \varphi_n(z)\rvert <\lvert \varphi_n^*(z)\rvert \quad \text{ and }\quad \lvert \Phi_n(z)\rvert <\lvert \Phi_n^*(z)\rvert, \quad z\in \mathbb{D}.
        \end{equation*}
        \item If $\varphi_0, \varphi_1, \ldots, \varphi_n$ exist and are unique, then they satisfy the following Christoffel-Darboux identity
        \begin{equation*}
            \sum_{j=0}^{n-1} \varphi_j(z)\overline{\varphi_j(w)}=\frac{\varphi_n^*(z)\overline{\varphi_n^*(w)}-\varphi_n(z)\overline{\varphi_n(w)}}{1-\overline{w}z},
        \end{equation*}
        for all $\overline{w}z\neq 1$. For $\overline{w}z=1$, the right-hand side of this identity is understood as a limit.
    \end{enumerate}
\end{lemma}

\begin{equation}\label{eq:choice-of-polynomial-p-as-orthonormal}
    P(z)=\varphi_{n+1}^*(z;\d\mu),    
\end{equation}

\begin{equation}\label{eq:choice-of-polynomial-p}
    P(z)=(1+i)\varphi_{n}^*(z;\d\mu) + i z \varphi_{n}(z;\d\mu).
\end{equation}

\begin{equation}\label{eq:reproducing-kernel-formula}
    K_n(w, z)=\frac{P(z) \overline{P(w)}-P^*(z) \overline{P^*(w)}}{1-\overline{w}z},
\end{equation}

\begin{lemma}\label{lem:quadrature-formula-for-laurent-polynomials}
Let $\mu$ be a probability measure on $\dD$ and $W\in \Gamma_n$ be a Laurent polynomial. Let $P$ be given by \eqref{eq:choice-of-polynomial-p-as-orthonormal} or by \eqref{eq:choice-of-polynomial-p}. For each fixed $\theta\in \R$ we have 
\begin{equation}\label{eq:quadrature formula for W}
    \int_{\dD} W(z)\, \d\mu(z)=\int_{\dD}\frac{W(z)}{|P(z)|^2}\,\d t=\sum_{T_\theta(\xi)=0}\frac{W(\xi)}{K_{n}(\xi, \xi)},
\end{equation}
where $z=e^{2\pi i t}$ and $T_\theta(z)=e^{i\theta}P(z)-e^{-i\theta}P^*(z)$.
\end{lemma}

\begin{proof}
Recall the choice of $P(z)$ from \eqref{eq:choice-of-polynomial-p-as-orthonormal} and \eqref{eq:choice-of-polynomial-p}. Using the representation \eqref{eq:reproducing-kernel-formula} of the reproducing kernel $K_n$ of the space $\mathcal{H}_n(P)$ and the Christoffel-Darboux identity from Lemma \ref{lem:properties of orthogonal polynomials} we obtain that 
\begin{equation}\label{eq:christoffel-darboux-identity-of-K}
    K_n(w, z)=\sum_{j=0}^n \varphi_j(z;\d\mu)\overline{\varphi_j(w; \d\mu)}.
\end{equation}

Fix $\theta\in \R$ and consider the polynomial $T_\theta(z)$ whose zeros are denoted by $\xi_1, \xi_2, \ldots, \xi_{n+1}$. Now consider the set of polynomials given by 
\begin{equation*}
    q_j(z)=\frac{K_n(\xi_j, z)}{\sqrt{K_n(\xi_j, \xi_j)}}, \quad 1\leq j\leq n+1.
\end{equation*} 
Recall the fact that $\mathcal{H}_n(P)=\mathcal{H}_n(e^{i\theta}P)$ and the formula \eqref{eq:reproducing-kernel-with-A-and-B}. Since $\xi_1, \xi_2, \ldots, \xi_{n+1}$ are distinct zeros of the polynomial $T_\theta$ we know the orthogonality of $K_n(\xi_j, \cdot)$ and $K_n(\xi_k, \cdot)$ in $\mathcal{H}_n(P)$.  Thus, we know $K_n(\xi_j, \xi_k)=0$, for all $j\neq k$ and therefore $q_1, q_2, \ldots, q_{n+1}$ is a basis of the space $\mathcal{P}_n$. Also, note that using \eqref{eq:christoffel-darboux-identity-of-K} we get 
\begin{equation*}
    \begin{split}
        \langle q_j, q_k\rangle_{L^2(\dD, \d\mu)}&=\frac{\langle K_n(\xi_j, \cdot), K_n(\xi_k, \cdot)\rangle_{L^2(\dD, \d\mu)}}{\sqrt{K_n(\xi_j, \xi_j)}\sqrt{ K_n(\xi_k, \xi_k)}}\\
        &=\frac{1}{\sqrt{K_n(\xi_j, \xi_j) K_n(\xi_k, \xi_k)}}\sum_{0\leq r, s\leq n}\overline{\varphi_r(\xi_j)}\varphi_s(\xi_k)\; \langle \varphi_r, \varphi_s\rangle_{L^2(\dD, \d\mu)}\\
        &=\frac{K_n(\xi_j, \xi_k)}{\sqrt{K_n(\xi_j, \xi_j) K_n(\xi_k, \xi_k)}}=\langle q_j, q_k\rangle_{\mathcal{H}_n(P)}.
    \end{split}
\end{equation*}
Since the $q_j$'s form a basis of the space $\P_n$, we conclude that the inner products $\langle \cdot , \cdot \rangle_{L^2(\dD, \d\mu)}$  and $\langle \cdot, \cdot \rangle_{\mathcal{H}_n(P)}$ are the same for elements of $\P_n$. 
\medskip

Now we come to the statement of the lemma. Because of the linearity, it suffices to prove the statement for $W_j(z)=z^j$, with $-n\leq j\leq n$.  First, consider the case $0\leq j\leq n$. We have 
\begin{equation*}
    \int_{\dD} z^j\, \d\mu(z)=\langle W_j, 1\rangle_{L^2(\dD, \mu)}=\langle W_j, 1\rangle_{\mathcal{H}_n(P)}=\int_{\dD}\frac{z^j}{|P(z)|^2}\, \d t.
\end{equation*}
Moreover, using the identity $\langle q_j, q_k\rangle=\delta_{j,k}$, we can write $1=\sum_{k=1}^{n+1} \langle 1, q_k\rangle q_k(z)$ and 
\begin{equation*}
    \langle W_j, 1\rangle_{\mathcal{H}_n(P)}=\sum_{k=1}^{n+1}\overline{\langle 1, q_k\rangle}\langle W_j, q_k\rangle=\sum_{k=1}^{n+1}\frac{W_j(\xi_k)}{K_n(\xi_k, \xi_k)},
\end{equation*}
where in the last identity we used the fact that $\langle Q, q_k\rangle=Q(\xi_k)/\sqrt{K_n(\xi_k, \xi_k)}$ for all polynomials $Q\in \P_n=\mathcal{H}_n(P)$.

For $-n\leq j\leq -1$, we write
\begin{equation*}
    \int_{\dD} z^j\, \d\mu(z)=\langle 1, W_{-j}\rangle=\int_{\dD} \frac{\overline{W_{-j}(z)}}{|P(z)|^2}\,\d t=\int_{\dD}\frac{z^j}{|P(z)|^2}\,\d t.
\end{equation*}
Similarly, we can conclude the final identity using the fact that $\overline{\xi_k^{-j}}=\xi_k^{j}$. 
\end{proof}

\end{document}
